\chapter{The contest playbook}\label{ch:playbook}
\index{contest strategy}\index{stress testing}\index{upsolving}\index{scoreboard}\index{time allocation}\index{practice ladder}

Everything before this chapter taught you to solve problems.
This last one is about the machine that wraps around solving:
the half-hour before the clock starts, the scoreboard telling
you where the points are, the decision to abandon a problem,
the brute-force oracle that catches the bug your tests missed,
and the week after, when the contest actually teaches you
something. Technique gets you to a score; process gets you to a
rating. After this chapter you can run a contest like an
engineer: a plan going in, telemetry during, and a triage
afterwards.

% from-zero:begin
\begin{rulebox}[Before you start]
All of the book up to Chapter~17 as a reference, but nothing new is required: this chapter is about \emph{process}. \emph{Recap:} the algorithms are the easy part of a contest; losing points to unread statements, untested code and badly managed time is what keeps most contestants below their real level.
\end{rulebox}
\begin{wins}[After this chapter you can]
\begin{itemize}
\item prepare a template, notes and environment before a contest starts;
\item allocate time among problems and know when to abandon one;
\item find a failing case by stress testing against a brute force;
\item run an upsolving routine that actually raises your rating.
\end{itemize}
\end{wins}
\begin{principle}[How to think about this]
Strong contestants are not the ones who know the most algorithms; they are the ones who waste the least time. The way of thinking is to treat a contest as a project: read everything first, solve the sure things, test before submitting, and learn from each mistake afterwards. After the contest, solving the problems you could not do is where most of the improvement happens.
\end{principle}
% from-zero:end

\section{Before the contest}\label{sec:before}
\index{pre-contest routine}
The contest is decided before it starts more often than people
admit. A fixed pre-contest routine removes the one-shot
decisions that cost real minutes under pressure. Thirty minutes
out: environment ready (compiler flags, template file, a
\texttt{diff}-able checker), a fresh physical notebook, water,
phone off. Ten minutes out: re-read the rules that cost points
--- submission limits, penalty per wrong try, whether partial
credit exists (IOI-style subtasks change the whole strategy;
see \S\ref{sec:styles}). One minute out: nothing. Starting
cold and calm beats starting warm and frantic.

\begin{example}[The ten-minute setup]
Write your setup down once and reuse it verbatim --- ours, used
throughout this book: template file loaded,
\texttt{g++ -O2 -std=c++17 -Wall} alias checked on a trivial
compile, \texttt{stress.sh} reachable, scratch page split into
``idea / invariant / test'' columns. It takes under ten
minutes and removes the most common minute-zero stalls.
\end{example}
\begin{solution}
The deliverable is a checklist, not a contest performance ---
its entire purpose is to be executed without thinking. Real
output of \texttt{stress.sh --help}-style first run on a fresh
box:
\begin{verbatim}
MISMATCH seed=15  brute=-15  fast=0
input:
1
-15
\end{verbatim}
--- proving the harness itself works before you need it.
\end{solution}

\section{Reading the scoreboard}\label{sec:board}
\index{scoreboard!reading}
The scoreboard is live information, not decoration. Three
numbers matter per problem: how many teams solved it
(difficulty proxy), how many \emph{tried} it (read effort
proxy), and your own attempts (sunk cost to resist). Solving
order follows the solved-count, not the problem letters --- the
crowd has already ranked the set for you
(Figure~\ref{fig:board}). ICPC rules make this sharper: penalty
minutes accumulate on wrong submissions to \emph{solved}
problems, so a speculative submit on a cold problem is cheap
(zero penalty if never solved) while a third wrong guess on a
problem you'll solve is expensive.
\begin{figure}[ht]\centering
\begin{tikzpicture}[font=\tiny,scale=0.72]
\draw[rulegray] (0,0) rectangle (10.6,3.4);
\foreach \y in {0.6,1.2,1.8,2.4,3.0} \draw[rulegray] (0,\y)--(10.6,\y);
\foreach \x/\t in {0.0/Rank,1.4/Team,4.0/A,5.0/B,6.0/C,7.0/D,8.0/E,9.0/Pen.}
  \node[anchor=west,font=\tiny\bfseries] at (\x+0.12,3.1) {\t};
\foreach \y/\rk/\tm/\a/\b/\c/\d/\e/\p in {
  2.7/1/byte girls/+2/+/{}/{}/+1/89:52,
  2.1/2/night owls/+/+1/{}/{}/-2/143:20,
  1.5/3/you/{}/{}/-3/{}/{}/---,
  0.9/4/segfaults/-1/{}/{}/{}/{}/---,
  0.3/5/newbies/{}/{}/{}/{}/{}/---}
  {\node[anchor=west] at (0.12,\y) {\rk};
   \node[anchor=west] at (1.52,\y) {\tm};
   \node at (4.5,\y) {\a}; \node at (5.5,\y) {\b}; \node at (6.5,\y) {\c};
   \node at (7.5,\y) {\d}; \node at (8.5,\y) {\e}; \node[anchor=east] at (10.5,\y) {\p};}
\node[align=left,anchor=west] at (10.9,1.7) {Read it as:\\$\bullet$ +: solved, --: failed tries\\$\bullet$ +2 = solved on 3rd try\\$\bullet$ penalty = time + 20m/try\\$\bullet$ unsolved problems with\\most tries = where effort goes};
\end{tikzpicture}

\caption{A scoreboard five minutes in. Notice what each cell
tells you: \texttt{+}/\texttt{-} counts reveal where attempts
are going, and the top rows' unsolved letters are the de-facto
hard end of the set.}
\label{fig:board}
\end{figure}

\begin{example}[Five-minute scoreboard read]
Snapshot: 4 of 7 problems solved by someone; B has 40
submissions and 30 solves; E has 3 tries, 0 solves. You have
solved nothing yet.
\textbf{Strategy.} B is the intended entry --- a $3/4$ solve
rate at volume means readable and mechanical. Skip A/C/D's
lower traffic for now; E is untouched for a reason. Your order:
B, then whichever of A/C shows the next-solve first (fresh
telemetry), E last or never.
\end{example}
\begin{solution}
Order set in under a minute, and it updates: if A gets the next
solve before C, A moves ahead. The scoreboard re-ranks the set
every few minutes --- check it whenever you finish a problem or
go quiet for ten minutes.
\end{solution}

\begin{exercise}[easy] On the same snapshot, argue why ``start
with E because nobody solved it, glory awaits'' is usually
wrong --- and name the one situation (your specialty topic)
where it is right.\end{exercise}

\section{Time allocation and switching}\label{sec:time}
\index{switching problems}\index{time cap}
A four-hour contest is a budget: survey ($\sim$20 min),
solve-in-rank-order, per-problem caps, reserve the last
25 minutes for submit-or-stress only
(Figure~\ref{fig:time}). The per-problem cap is the rule
everyone agrees with and nobody obeys: \textbf{25 minutes of
no new idea means switch}. Not 35, not ``one more thing to
try.'' The trigger is not elapsed time but \emph{progress}
--- a new formulation, a shrinking search space, a passing
test. Staring blankly at the statement for the fifth time is
the canonical no-progress signal (Figure~\ref{fig:switch}).
\begin{figure}[ht]\centering
\begin{tikzpicture}[font=\tiny,scale=0.78,x=1.35cm]
\foreach \x in {0,0.5,...,4} \draw[rulegray!60] (\x,0)--(\x,2.6);
\foreach \x/\h in {0/0h,1/1h,2/2h,3/3h,4/4h} \node[below] at (\x,0) {\h};
\node[left] at (0,2.35) {plan};
\node[left] at (0,1.55) {read all};
\node[left] at (0,0.95) {solve order};
\node[left] at (0,0.35) {buffer};
\fill[accentsoft] (0,2.15) rectangle (0.33,2.55);
\fill[good] (0.33,1.35) rectangle (1.0,1.75);
\fill[accent]  (1.0,1.35) rectangle (2.2,1.75);
\fill[amber]   (2.2,1.35) rectangle (3.6,1.75);
\fill[rulegray] (3.6,2.15) rectangle (4.0,2.55);
\node[above=2pt] at (0.16,2.55) {survey};
\node[align=left,anchor=west] at (4.15,1.5) {rules:\\$\bullet$ 20 min survey first\\$\bullet$ easiest first,\\
no skipping twice\\$\bullet$ hard cap per\\problem (90 min)\\$\bullet$ last 25 min:\\submit-or-stress only};
\end{tikzpicture}

\caption{A 4-hour budget that survives contact. Notice the
plan bar vs.\ the work bars: the buffer is scheduled, not
stolen; the last block is for submitting and stress-checking
only, never for starting something new.}
\label{fig:time}
\end{figure}
\begin{figure}[ht]\centering
\begin{tikzpicture}[font=\tiny,scale=0.78]
\node[draw=accent,rounded corners=2pt,align=center] (a) at (0,0) {stuck \textgreater cap?\\(rule: 25 min no progress)};
\node[draw=slate,rounded corners=2pt,align=center] (b) at (4.3,1.3) {leave it: restate problem\\on scratch, list 3 approaches};
\node[draw=slate,rounded corners=2pt,align=center] (c) at (4.3,-1.3) {switch to next readable\\problem};
\node[draw=good,rounded corners=2pt,align=center] (d) at (8.6,0) {return only with a\\new idea in hand};
\draw[->,slate] (a)--node[above,sloped]{idea exists}(b);
\draw[->,slate] (a)--node[below,sloped]{staring blank}(c);
\draw[->,good,dashed] (b.east)--(d.north);
\draw[->,good,dashed] (c.east)--(d.south);
\node[align=left,anchor=west] at (9.4,0) {the real rule:\\switching is free,\\re-reading is what\\costs time};
\end{tikzpicture}

\caption{The switching decision. Notice it asks ``do you have
a new idea,'' not ``have you spent long enough'' --- an
unexplored approach resets the clock, a blank stare does not.}
\label{fig:switch}
\end{figure}

\begin{example}[A contest log, annotated]
Minute 0--18: survey all 7, sketch data ranges. 18--55: solve
B (impl, first try). 55--80: D attempt, two wrong tries on
samples, no new idea at minute 80 $\to$ switch. 80--110: solve
C. 110--135: back on D with the trick seen while doing C ---
solved. 135--200: stress A against brute; found and fixed an
edge case, submitted. 200--240: writeup of E abandoned at 15\%
confidence --- final reserve held.
\textbf{Strategy.} The log shows the two habits that pay:
switching on a progress trigger (not a clock), and returning
with ammunition rather than stubbornness.
\end{example}
\begin{solution}
Three solves plus a fourth saved by the stress reserve. The
same contest played ``in order, no switches'' ends at two.
\end{solution}

\begin{exercise}[medium] Write your own caps: survey length,
per-problem cap, switch trigger, reserve length --- for a 2-hour
CF Div.~2 \emph{and} a 5-hour ICPC regional. Defend why the
numbers differ.\end{exercise}

\section{Stress testing}\label{sec:stress}
\index{stress testing!harness}\index{brute-force oracle}
The strongest debugging weapon in this book is two programs
and a loop: a \emph{brute force} so simple it is obviously
correct, your fast submission, a generator of small random
inputs, and a comparator that stops on the first difference
(Figure~\ref{fig:stress}). Small inputs are the point: bugs
live in edge cases (empty, singleton, all-equal, negative,
overflow-adjacent), and small cases let you run thousands of
rounds per second while still being checkable by hand when one
fails.
\begin{figure}[ht]\centering
\begin{tikzpicture}[font=\tiny,node distance=4mm and 8mm,scale=0.78]
\node[draw=accent,rounded corners=2pt,minimum width=2.0cm,minimum height=0.7cm,align=center] (gen) {generator\\(random seed)};
\node[draw=slate,rounded corners=2pt,minimum width=2.0cm,minimum height=0.7cm,align=center,right=of gen] (inp) {one small\\input};
\node[draw=slate,rounded corners=2pt,align=center,right=of inp,yshift=7mm] (fast) {fast (submission)};
\node[draw=slate,rounded corners=2pt,align=center,right=of inp,yshift=-7mm] (brute) {brute (obviously right)};
\node[draw=accent,rounded corners=2pt,align=center,right=of fast,xshift=2mm,yshift=-7mm] (cmp) {outputs\\equal?};
\node[draw=good,rounded corners=2pt,right=of cmp] (pass) {next seed};
\node[draw=amber,rounded corners=2pt,align=center,below=4mm of pass] (fail) {STOP: print input\\+ both answers};
\draw[->,slate] (gen)--(inp);
\draw[->,slate] (inp)--(fast.west);
\draw[->,slate] (inp)--(brute.west);
\draw[->,slate] (fast)--(cmp);
\draw[->,slate] (brute)--(cmp);
\draw[->,good] (cmp)--node[above]{same}(pass);
\draw[->,amber] (cmp)|-(fail) node[pos=0.25,below]{diff};
\draw[->,rulegray,dashed] (fail.west)-|([xshift=-3mm]gen.south west)|-(gen);
\node[below=1mm of fail,align=center] {then shrink the\\failing input by hand};
\end{tikzpicture}

\caption{The harness. Notice there is no ``is it right?''
step --- only ``do they agree?'' Correctness is borrowed from
the brute force; the fast program only has to match it.}
\label{fig:stress}
\end{figure}

\begin{example}[The harness, working]
Chapter 1's maximum-subarray problem: Kadane's against the
$O(n^3)$ triple loop, on 1--12 element arrays in $[-20,20]$.
We feed it \texttt{mss\_buggy}, a Kadane variant that lets the
empty segment in by starting \texttt{cur} and \texttt{best}
at zero.
\end{example}
\begin{solution}
\lstinputlisting{code/gen_mss.cpp}
\lstinputlisting{code/mss_buggy.cpp}
\lstinputlisting{code/stress.sh}
Real output:
\begin{verbatim}
MISMATCH seed=15  brute=-15  fast=0
input:
1
-15
\end{verbatim}
Fifteen seeds, and the harness hands you a one-element failing
case: all-negative input, where the buggy version's empty
segment answers $0$ instead of $-15$. Rerun on the corrected
\texttt{mss\_linear}:
\begin{verbatim}
stress: 500/500 rounds identical (mss_linear vs mss_brute)
\end{verbatim}
\end{solution}

The counterexample it gives you is usually already small; when
it is not, \emph{shrink} it: delete elements while the outputs
still differ. A failing case you cannot explain is a failing
case you cannot fix --- the shrink loop is where the
understanding lives.

\begin{example}[When to reach for it]
Rules of thumb: any greedy that ``looks right'' (ch.~3), any
DP with boundary states (chs.~11--12), anything with a
special case you added last (chs.~16--17 all hit this ---
\texttt{seg\_int}'s endpoint-touch and \texttt{digitdp}'s
adjacent rule were both harness-caught). If the problem has
no cheap oracle, write the oracle: a slower correct solution
exists for almost every problem in this book.
\end{example}
\begin{solution}
The cost is real --- 20--40 minutes to build the pair. The
return is a wrong-answer converted to a failing input, which
is the difference between debugging and guessing.
\end{solution}

\begin{exercise}[medium] Build the pair for ch.~12's
\emph{Counting Numbers}: the DP solver plus a brute
\texttt{for(int x=a;x<=b;x++)} checker on small ranges.
What input class did your own solver's bug live in?\end{exercise}

\section{Upsolving}\label{sec:upsolve}
\index{upsolving!discipline}
The contest ends; the learning starts. Every problem you
failed or skipped goes into one of three buckets within the
week: knew-it-but-botched (rewrite cold, no reference), idea
was close (one editorial line, then 30 more minutes alone),
unknown technique (editorial + two more problems using it).
Figure~\ref{fig:upsolve} makes it a pipeline with one exit:
an accepted submission \emph{and} a one-line lesson in your
log. ``I knew that one'' is not a lesson; ``break ties by the
right endpoint, not the left'' is.
\begin{figure}[ht]\centering
\begin{tikzpicture}[font=\tiny,scale=0.78]
\node[draw=accent,rounded corners=2pt,align=center] (u) at (0,0) {every unsolved+failed\\problem, same week};
\node[draw=slate,rounded corners=2pt,align=center] (a) at (4.4,1.5) {knew the idea,\\botched impl:\\rewrite cold};
\node[draw=slate,rounded corners=2pt,align=center] (b) at (4.4,0) {idea was close:\\read 1 hint line,\\retry 30 min};
\node[draw=slate,rounded corners=2pt,align=center] (c) at (4.4,-1.5) {unknown technique:\\read editorial,\\solve 2 more of it};
\draw[->,slate] (u)--(a.west); \draw[->,slate] (u)--(b.west); \draw[->,slate] (u)--(c.west);
\node[draw=good,rounded corners=2pt,align=center] (d) at (8.9,0) {submit accepted,\\file 1-line lesson\\in your log};
\draw[->,good,dashed] (a.east)--(d.north); \draw[->,good,dashed] (b.east)--(d); \draw[->,good,dashed] (c.east)--(d.south);
\node[below=4mm of c,align=left] {upsolve $=$ the contest's real teaching;\\skipping it wastes the entry fee};
\end{tikzpicture}

\caption{Upsolve triage. Notice each bucket has a different
prescription --- reading the editorial is the \emph{last}
bucket's tool, and applying it to the first two is how a
contest teaches nothing.}
\label{fig:upsolve}
\end{figure}

\begin{example}[Triage after a Div.~2]
Post-contest pile: B (WA on test 4), D (skipped), E (no idea).
B: your comparator forgot equal keys --- bucket 1, rewrite
cold, accepted in 15 min, log: ``sort ties break on index.''
D: one hint line (``think backwards'') $\to$ solved in 40 min,
bucket 2. E: editorial says centroid decomposition --- bucket
3: implement it, then find two centroid problems and log the
pattern.
\end{example}
\begin{solution}
Three problems, three different actions, three log lines.
That is one contest converted to durable skill; skip the week
and it was only a number.
\end{solution}

\begin{exercise}[easy] Take your last practice set and assign
every unsolved item to a bucket. If everything lands in
bucket 3, your contest level is too high --- explain
why.\end{exercise}

\section{The practice ladder and contest styles}\label{sec:styles}
\index{practice ladder}\index{ICPC}\index{IOI}
Rating climbs by repetition at the right altitude
(Figure~\ref{fig:ladder}): solve where your success rate is
40--60\%, move up when it passes 60\%, move down after three
consecutive failures, and draw from named banks (CSES for
fundamentals, USACO Guide lists for technique depth, CF
div.~2 A--C for speed). Adjust the machinery per style:
\textbf{ICPC} is teams, penalty minutes, and scoreboard
meta-play --- pick off the solved-count leaders and guard
your penalty. \textbf{IOI} is solo, subtask-scored, five
hours --- partial credit means a brute-force subtask taken
early is worth real points \emph{and} doubles as your stress
oracle for the full score. \textbf{CF} sits between: short
rounds, systests, and pretest gamesmanship --- never trust a
pretest-pass as final.
\begin{figure}[ht]\centering
\begin{tikzpicture}[font=\tiny,scale=0.78]
\draw[->,rulegray] (0,0)--(11,0) node[below left] {CF rating};
\foreach \x/\r/\t in {0.4/800/impl+math,1.8/1100/greedy+sim,3.4/1300/dp+graphs lite,5.2/1600/full toolkit,7.4/1900/adv data structs,9.4/2100+/creative}
  {\fill[accentsoft] (\x,0) rectangle (\x+1.2,0.5+0.12*\x);
   \node at (\x+0.6,-0.4) {\r};
   \node[align=center,below] at (\x+0.6,-0.7) {\t};}
\node[align=left,anchor=west] at (0,3.2) {ladder rules:\\$\bullet$ solve rate $>60\%$: move up\\$\bullet$ 3+ fails in a row: move down, upsolve the misses\\$\bullet$ bank named sets: CSES, USACO guide lists, CF div2 A--C};
\end{tikzpicture}

\caption{The 800$\to$2100 ladder. Notice the height is solve
rate, not volume: the ladder moves you, not the other way
around.}
\label{fig:ladder}
\end{figure}

\begin{example}[A twelve-week plan]
Weeks 1--4: CSES Introductory set (Weird Algorithm through
Two Knights) --- unranked, fundamentals only, target 20+
solves. Weeks 5--8: CF div.~2 A--B live, upsolve every C you
miss; weekly contest regardless of result. Weeks 9--12:
CSES Sorting/DP overlap + div.~2 C as the tracked metric.
Target: consistent div.~2 A+B, occasional C --- rating
$=$ consistency, not peak.
\end{example}
\begin{solution}
The plan's only hard rule is the upsolve loop from
\S\ref{sec:upsolve}; everything else is scheduling.
\end{solution}

\section{Common mistakes}
\index{contest strategy!common mistakes}
\begin{itemize}
\item \textbf{Solving in statement order.} Problem letters are
not difficulty ranks --- the scoreboard's solve counts are.
\item \textbf{The sunk-cost stall.} ``I've spent an hour''
argues for leaving, not staying; the cap triggers on
progress, not investment.
\item \textbf{Guessing a fix and resubmitting.} ICPC penalty
minutes and CF failed-test locks both tax blind retries ---
stress first, submit once.
\item \textbf{No brute force, ever.} ``Too slow for the
judge'' is the feature: the oracle only has to be right.
\item \textbf{Upsolving by reading the editorial end-to-end.}
It converts an unknown-technique problem into a trivia fact
--- the bucket-3 treatment needs implementation, not
comprehension.
\item \textbf{Practicing only where you win.} A 90\% solve
rate is a flat ladder: the altitude is wrong, not the volume.
\end{itemize}

\section*{Summary}
\begin{tabularx}{\linewidth}{@{}lXl@{}}
\textbf{Moment} & \textbf{Move} & \textbf{Tool}\\
\midrule
before clock & fixed checklist, harness ready & \S\ref{sec:before}\\
first 20 min & survey all problems, sketch ranges & \S\ref{sec:board}\\
rank order & solve by solved-count, not letters & scoreboard\\
stuck & progress cap, not time cap & \S\ref{sec:time}\\
before submit & brute-vs-fast until identical & \texttt{stress.sh}\\
same week & triage fail/skip into 3 buckets & \S\ref{sec:upsolve}\\
between rounds & 40--60\% altitude, named banks & \S\ref{sec:styles}\\
\end{tabularx}

\section*{Exercises}
\begin{exercise}[easy] Write your pre-contest checklist (10 lines max), then time yourself executing it. \end{exercise}
\begin{exercise}[easy] For a contest you already played, reconstruct the solve order the scoreboard was telling you at minute 30 vs.\ the order you actually used. \end{exercise}
\begin{exercise}[easy] Implement \texttt{gen}/\texttt{brute}/\texttt{fast} for ch.~3's coin-change greedy vs.\ its DP --- the harness should find the canonical counterexample. \end{exercise}
\begin{exercise}[easy] List the three smallest inputs that break \texttt{mss\_buggy}-class logic (empty, singleton negative, all-negative) and explain why each is the \emph{minimal} witness. \end{exercise}
\begin{exercise}[medium] Time-cap experiment: run a 2-hour set twice --- once with a hard 25-minute switch rule, once ad-lib --- and compare solves and stress. \end{exercise}
\begin{exercise}[medium] Extend \texttt{stress.sh} with input \emph{shrinking}: on MISMATCH, try deleting each element while outputs still differ. Report how small your failing cases get. \end{exercise}
\begin{exercise}[medium] Upsolve your last 10 fails with the 3-bucket triage; log one line each. What fraction were bucket 1 --- and what does that say about your bottleneck? \end{exercise}
\begin{exercise}[medium] Build a personal ``first-15-minutes'' routine: what you read, what you compute per problem, what you write down. Test it on a virtual contest. \end{exercise}
\begin{exercise}[medium] ICPC vs.\ IOI scoring: take one problem and describe the optimal attempt plan under each --- where does brute-forcing a subtask sit in each plan? \end{exercise}
\begin{exercise}[medium] Design a generator that adversarially covers your bug history: for each past WA, encode the input class that caught it. Run it as regression after every new solver. \end{exercise}
\begin{exercise}[hard] Auto-verify a whole chapter: wire \texttt{stress.sh}-style oracles for five ch.~11--17 solvers, run 200 rounds each, and record which solver needed the largest generator space to stay green. \end{exercise}
\begin{exercise}[hard] Model the scoreboard as data: given per-problem solve counts over time, write the re-ranking rule \S\ref{sec:board} uses informally and simulate its chosen order on a real contest's history. \end{exercise}
\begin{exercise}[hard] Prove to yourself the ladder works: track solve-rate vs.\ CF rating over 30 problems and test the 60\% rule against your own trajectory (a spreadsheet is fine). \end{exercise}
\begin{exercise}[hard] Write a ``contest minute'' postmortem template --- timeline, decision points, what the scoreboard said vs.\ what you did --- and use it on your next three contests; report what changed. \end{exercise}
\begin{exercise}[hard] For a five-hour IOI-style set, write the attempt plan that maximizes expected partial credit given per-subtask point values; show where full-score attempts are worth their risk. \end{exercise}

\chapternotes

\section*{Further reading}
\begin{itemize}
\item USACO Guide, \emph{Practicing CP} --- the training
method this chapter's ladder compresses; del Nor\'e's
what-to-practice framework.
\newline\url{https://usaco.guide/general/practicing}
\item A.~Laaksonen, \emph{Competitive Programmer's Handbook}
--- its advice chapters on contest technique and debugging.
\newline\url{https://cses.fi/book/book.pdf}
\item Topcoder Thrive, \emph{How To Find a Solution} ---
classic problem-attack walkthroughs, same lineage as the
survey-first habit. (Thrive articles $\to$ title above, under:)
\newline\url{https://www.topcoder.com/}
\item ICPC, \emph{World Finals Rules} --- the penalty model
\S\ref{sec:board} reads from; worth reading once, primary
source.
\newline\url{https://icpc.global/worldfinals/rules}
\item Kenkoooo AtCoder Problems --- the other major problem
bank with rating tables; the ladder's alternate track.
\newline\url{https://kenkoooo.com/atcoder/}
\item cp-algorithms --- the reference for every technique an
upsolve bucket-3 will send you to.
\newline\url{https://cp-algorithms.com/}
\item infoarena \emph{evaluator} notes --- how judges and
partial scoring work, the IOI-side view.
\newline\url{https://infoarena.ro/blog/evaluator}
\item Codeforces problem set --- the bank itself
[unverified from this session: CF serves 403 to curl].
\newline\url{https://codeforces.com/problemset}
\end{itemize}

\begin{rulebox}[Watch instead]
\begin{itemize}
\item \textbf{Errichto --- ``How to test your solution in
Competitive Programming, on Linux?''}, 26m56s,
\url{https://www.youtube.com/watch?v=JXTVOyQpSGM} --- covers
\S\ref{sec:stress}: the brute-vs-fast methodology built live.
Watch this first. (Verified 2026-10-11.)
\item \textbf{Colin Galen --- ``Competitive Programming: How
to Learn from Contests''}, 20m58s,
\url{https://www.youtube.com/watch?v=E-aylp6MZnM} --- covers
\S\ref{sec:upsolve}--\ref{sec:styles}: post-contest analysis
and rating mechanics. (Verified 2026-10-11.)
\item \textbf{Errichto --- ``How to start Competitive
Programming? For beginners!''}, 9m43s,
\url{https://www.youtube.com/watch?v=xAeiXy8-9Y8} --- covers
\S\ref{sec:before}--\ref{sec:board}: the survey and setup
mentality. (Verified 2026-10-11.)
\end{itemize}
\emph{Notation note:} Galen uses ``rating'' for our ``ladder
altitude''; Errichto's Linux shell maps one-to-one onto
\texttt{stress.sh}.
\end{rulebox}
The harness methodology follows the standard
generator--oracle--comparator pattern as practiced on this
book's own code; the upsolve and ladder framing follows the
USACO Guide's practice guidance~\cite{usacoguide}; scoring
models from the ICPC rules~\cite{icpcrules} and the CSES
problem set structure~\cite{cses}.
